% ---------------------------------------------------------------------------
% Abstract / Kurzfassung
% The English abstract is wrapped in an English-language block so that line
% breaking uses the correct *English* hyphenation patterns; the surrounding
% document keeps German hyphenation.
% ---------------------------------------------------------------------------
\chapter*{Abstract}
\addcontentsline{toc}{chapter}{Abstract}

\begin{otherlanguage}{english}
\noindent
This bachelor thesis replicates the study ``Greater local supply of language courses improves refugees' labor market integration'' by Kanas and Kosyakova and extends it by treating asylum status as a moderator. The analyses draw on the IAB-BAMF-SOEP Survey of Refugees, waves 2016 to 2019, combined with county-level data on the supply of integration courses; they were run through SOEPremote, the remote execution system of DIW Berlin. The replication reproduces the sample selection exactly (5,467 person-years from 2,730 persons) and confirms the central result: a one standard deviation higher course supply in the county of initial assignment is associated with a 2.98 percentage point higher probability of employment (published: 2.97).

The extension builds on Esser's action-theoretical model of language acquisition and asks whether refugees with a recognised asylum claim, who have a secure prospect of staying and unrestricted access to courses, benefit more from the local supply than refugees whose claim is pending or rejected. This is tested not only for employment but, above all, for the mechanisms of integration course participation, course completion, language certificates and German language proficiency. None of the joint interaction tests is significant; course participation comes closest (p = 0.053), with point estimates ordered as expected. The results do not support the hypothesis, but given the group sizes they do not rule out smaller differences either.
\end{otherlanguage}

\bigskip
\begin{otherlanguage}{ngerman}
\noindent\textbf{Kurzfassung.}\enspace

Diese Bachelorarbeit reproduziert die Studie „Greater local supply of language courses improves refugees' labor market integration“ von Kanas und Kosyakova und erweitert sie um den Asylstatus als Moderator. Grundlage sind die IAB-BAMF-SOEP-Befragung von Geflüchteten in den Wellen 2016 bis 2019 sowie Kreisdaten zum Angebot an Integrationskursen; die Analysen liefen über das Ferndatenverarbeitungssystem SOEPremote des DIW Berlin. Die vorliegende Arbeit reproduziert die Stichprobenselektion exakt (5.467 Personenjahre von 2.730 Personen) und bestätigt das zentrale Ergebnis: Ein um eine Standardabweichung höheres Kursangebot im Zuweisungskreis geht mit einer um 2,98 Prozentpunkte höheren Beschäftigungswahrscheinlichkeit einher (publiziert: 2,97). 

Die Erweiterung prüft auf Basis des handlungstheoretischen Spracherwerbsmodells von Esser, ob anerkannte Geflüchtete wegen sicherer Bleibeperspektive und uneingeschränktem Kurszugang stärker vom lokalen Angebot profitieren als Geflüchtete mit offenem oder abgelehntem Verfahren. Geprüft wird dies nicht nur für die Beschäftigung, sondern vor allem für die Mechanismen Integrationskursteilnahme, Integrationskursabschluss, Sprachzertifikat und Deutschkenntnisse. Kein gemeinsamer Interaktionstest ist signifikant; am nächsten kommt die Kursteilnahme (p = 0,053), bei der die Punktschätzungen in der erwarteten Reihenfolge liegen. Die Ergebnisse stützen die Hypothese nicht, schließen kleinere Unterschiede angesichts der Gruppengrößen aber auch nicht aus.

\end{otherlanguage}
